%% ma: L11-B11-0007
%% lop: 11
%% chuong: 4
%% bai: 11
%% dang: 11.11.4
%% loai: TL
%% muc: VD
%% dap_an: "-"
%% nguon: "de-cuong-giua-ky-1-toan-11-nam-2026-2027-truong-thpt-hoang-van-thu-ha-noi-2-8.pdf (M08, Đề cương giữa kì 1 Toán 11 THPT Hoàng Văn Thụ, Hà Nội), Đề minh họa, Phần IV Câu 5"
%% kien_thuc: "Bài 10 (giao tuyến, giao điểm), Bài 11 (định lí về giao tuyến song song), định lí Thalès (lớp 8)"
%% trang_thai: cho-duyet
%% ly_do: "Dạng mới đề xuất 11.11.4"
%% ghi_chu: "Đã kiểm bằng tọa độ: SQ = 3/2 SD, Q nằm ngoài đoạn SD; QS/QD = 3."
\begin{ex}
Cho hình chóp $S.ABCD$ có đáy $ABCD$ là hình bình hành. Gọi $M$, $N$ lần lượt là trung điểm của $SA$, $SB$. Điểm $P$ thuộc đoạn thẳng $BC$ sao cho $BP=3PC$.

a) Tìm giao tuyến của $(SAC)$ và $(SBD)$.

b) Gọi $Q$ là giao điểm của đường thẳng $SD$ và mặt phẳng $(MNP)$. Tính tỉ số $\dfrac{QS}{QD}$.
\loigiai{a) Gọi $O=AC\cap BD$. Khi đó $(SAC)\cap(SBD)=SO$.

b) $MN\parallel AB$ nên $(MNP)\cap(ABCD)$ là đường thẳng qua $P$ song song với $AB$, cắt $AD$ tại $R$ với $\dfrac{AR}{AD}=\dfrac{BP}{BC}=\dfrac34$. Do đó $(MNP)\cap(SAD)=MR$ và $Q=MR\cap SD$.

Trong $\triangle SAD$, kẻ $DT\parallel SA$ ($T\in MR$): $\dfrac{DT}{AM}=\dfrac{RD}{RA}=\dfrac13$, nên $DT=\dfrac13AM=\dfrac13SM$. Vì $DT\parallel SM$: $\dfrac{QS}{QD}=\dfrac{SM}{DT}=3$ ($Q$ nằm trên đường thẳng $SD$, ngoài đoạn $SD$, về phía $D$).}
\end{ex}
